% !TEX program = xelatex
% 吕博文 个人简历（中文版）— xelatex 编译
\documentclass[10pt,a4paper,fontset=none]{ctexart}

%==================== 字体 ====================
\usepackage{fontspec}
% 中文：系统自带 宋体 / 黑体（macOS）
\setCJKmainfont{Songti SC}[BoldFont={Heiti SC}]
\setCJKsansfont{Heiti SC}
\setCJKmonofont{Heiti SC}
% 西文：Times 风格衬线 + 无衬线标题（TeX Gyre，按文件名加载）
\setmainfont{texgyretermes}[
  Extension = .otf, UprightFont = *-regular, BoldFont = *-bold,
  ItalicFont = *-italic, BoldItalicFont = *-bolditalic]
\setsansfont{texgyreheros}[
  Extension = .otf, UprightFont = *-regular, BoldFont = *-bold,
  ItalicFont = *-italic, BoldItalicFont = *-bolditalic]

%==================== 版面 ====================
\usepackage[a4paper,top=1.1cm,bottom=1.0cm,left=1.5cm,right=1.5cm]{geometry}
\usepackage[dvipsnames,table]{xcolor}
\usepackage{titlesec}
\usepackage{enumitem}
\usepackage{graphicx}
\usepackage{tikz}
\usepackage{fontawesome5}
\usepackage[hidelinks]{hyperref}
\usepackage{ragged2e}

\definecolor{accent}{RGB}{31,78,214}      % 主题蓝
\definecolor{accentd}{RGB}{18,46,120}     % 深蓝
\definecolor{inkc}{RGB}{27,36,51}          % 正文墨色
\definecolor{mutedc}{RGB}{91,107,130}      % 次要灰
\definecolor{rulec}{RGB}{205,221,248}      % 浅色分隔线

\color{inkc}
\hypersetup{colorlinks=true, urlcolor=accent, linkcolor=accent}

\setlength{\parindent}{0pt}
\linespread{1.05}

%==================== 章节标题 ====================
\titleformat{\section}
  {\sffamily\large\bfseries\color{accentd}}
  {}{0pt}
  {}
  [{\vspace{2pt}\color{rulec}\titlerule[1.2pt]}]
\titlespacing*{\section}{0pt}{7pt}{4pt}

%==================== 小工具 ====================
% 一条经历的标题行：左标题（粗）+ 右时间（灰斜）
\newcommand{\entry}[2]{%
  \par\vspace{1pt}%
  {\sffamily\bfseries\color{inkc}#1}\hfill{\small\itshape\color{mutedc}#2}\par
  \vspace{0.5pt}}
% 副标题 / 角色行
\newcommand{\role}[1]{{\small\color{accentd}\sffamily #1}\par\vspace{1pt}}

\setlist[itemize]{leftmargin=1.4em, topsep=1pt, itemsep=1.5pt, parsep=0pt, label={\color{accent}\small$\bullet$}}

% 圆形头像（参数：图片、半径、宽度、竖向偏移）
\newcommand{\circphoto}[4]{%
  \begin{tikzpicture}
    \clip (0,0) circle (#2);
    \node at (0,#4) {\includegraphics[width=#3]{#1}};
    \draw[rulec, line width=1.2pt] (0,0) circle (#2);
  \end{tikzpicture}}

\begin{document}
\pagestyle{empty}

%==================== 抬头 ====================
\begin{minipage}[c]{0.72\textwidth}
  {\sffamily\fontsize{26}{30}\selectfont\bfseries\color{accentd}吕博文}\quad
  {\sffamily\large\color{mutedc}Bowen Lv}\\[4pt]
  {\sffamily\color{inkc}清华大学\,·\,计算机科学与技术\,·\,博士在读}\\[2pt]
  {\small\color{mutedc}研究方向：大语言模型（LLM）后训练、强化学习（RL）、GUI / Computer-Use Agent}\\[8pt]
  {\small
   \faEnvelope[regular]~\href{mailto:lvbowen1228@gmail.com}{lvbowen1228@gmail.com}\quad
   \faPhone*~18337283380\quad
   \faWeixin~18337283380\\[3pt]
   \faGlobe~\href{https://extreme1228.github.io/extreme1228/}{extreme1228.github.io/extreme1228}\quad
   \faGoogle~\href{https://scholar.google.com/citations?user=5JW6fqwAAAAJ&hl=zh-CN}{Google Scholar}}
\end{minipage}\hfill
\begin{minipage}[c]{0.24\textwidth}
  \raggedleft\circphoto{photo.jpg}{1.85cm}{3.7cm}{-0.55cm}
\end{minipage}

\vspace{2pt}

%==================== 教育背景 ====================
\section{教育背景}
\entry{清华大学\quad 计算机科学与技术系（KEG 知识工程实验室）}{2025 -- 至今}
\role{计算机科学与技术专业\;·\;博士研究生\;·\;导师：东昱晓\;老师（方向：LLM + 强化学习）}

\entry{同济大学\quad 电子与信息工程学院}{2021 -- 2025}
\role{计算机科学与技术专业（计算机拔尖实验班）\;·\;工学学士\;·\;专业排名 1 / 20}

%==================== 实习经历 ====================
\section{实习经历}
\entry{智谱 AI（Zhipu AI）\quad 大模型后训练 / GUI Agent RL\;算法实习生}{2025.02 -- 至今}
\role{主要负责大语言模型后训练与 GUI / Computer-Use Agent 的强化学习}
\begin{itemize}
  \item 深度参与 Agentic RL 训练基础设施建设：从零构建并重构出高度异步、训推分离的在线 RL 训练框架 \textbf{AgentRL}，集成多环境 Docker 化部署与高并发 rollout，支撑多任务多轮次（Multi-Task Multi-Turn）训练。
  \item 主导 Computer-Use Agent 的可验证数据合成与在线 RL 训练，提出 VeriGen 数据流水线、Frontier Sampling 采样与 Visual Context Segmentation 训练策略，在 OSWorld / ScienceBoard 上达到开源 SOTA（对应论文 \textbf{ScaleCUA}，第一作者）。
  \item 主导 \textbf{Auto-Post-Train（APT）} —— 一套 \textbf{AI for AI、模型自我进化（self-improvement）}的多智能体无人值守后训练系统：以 Conductor 统一调度 env / task / rollout / train 等多类 Worker，将\textbf{数据采集、环境适配、任务构建、rollout、质量审计到模型训练}的全流程完全自动化，让大模型在最少人工干预下持续自我迭代。
\end{itemize}

%==================== 科研成果 ====================
\section{科研成果}
% —— 着重突出 ScaleCUA ——
\entry{ScaleCUA: Scaling Computer Use Agents with Verifiable Task Synthesis and Efficient Online RL}{2026,\ 投稿中}
\role{Bowen Lv\;（第一作者）}
\begin{itemize}
  \item 提出统一框架 \textbf{ScaleCUA}，将可验证任务合成与高效在线 RL 结合，系统性破解 Computer-Use Agent 的数据稀缺与在线 RL 低效两大瓶颈；设计端到端数据流水线 \textbf{VeriGen}，借助共享 Docker 交互探针扩展到 100+ 并发 Worker，产出 24K+ 可验证任务与近 3K 高质量 RL 任务。
  \item 训练侧提出 \textbf{Frontier Sampling}（按模型能力前沿动态分配 rollout）与 \textbf{Visual Context Segmentation}（滑动窗口视觉上下文，训练提速 2.83$\times$）；以 Qwen3.5-9B 为基座，在 \textbf{OSWorld 68.7\%}、\textbf{ScienceBoard 54.0\%} 刷新开源 SOTA，超越 EvoCUA-32B、Kimi K2.5 等更大模型。
\end{itemize}

\entry{AgentRL: Scaling Agentic Reinforcement Learning with a Multi-Turn, Multi-Task Framework}{Preprint, 2025}
\role{Hanchen Zhang, Xiao Liu, \textbf{Bowen Lv}, et al.\quad\faExternalLink*~\href{https://arxiv.org/abs/2510.04206}{arXiv:2510.04206}}
\begin{itemize}
  \item 多任务、多轮次的智能体强化学习框架；RL 训练后平均性能在 AgentBench 上超越 Claude-4、GPT-5 等领先闭源 API。
\end{itemize}

\entry{KARL: Reinforcement Learning for LLM Agents on Multi-Turn Knowledge-Intensive Agentic Tasks}{ACL 2026 (Main)}
\role{Xueqiao Sun, Xiao Liu, \textbf{Bowen Lv}, et al.}
\begin{itemize}
  \item 面向多轮、知识密集型智能体任务的强化学习方法，被 ACL 2026 主会录用。
\end{itemize}

\entry{Focus on What Matters: Separated Models for Visual-Based RL Generalization}{NeurIPS 2024}
\role{Di Zhang, \textbf{Bowen Lv}, et al.\quad\faExternalLink*~\href{https://arxiv.org/abs/2410.10834}{arXiv:2410.10834}}
\begin{itemize}
  \item 以分离式模型提升视觉强化学习在复杂场景下的泛化能力，\textbf{第二作者}，发表于 NeurIPS 2024。
\end{itemize}

\entry{GLM-5V-Turbo: Toward a Native Foundation Model for Multimodal Agents}{Preprint, 2026}
\role{Wenyi Hong, Xiaotao Gu, \ldots, \textbf{Bowen Lv}, \ldots, GLM-V Team\quad\faExternalLink*~\href{https://arxiv.org/abs/2604.26752}{arXiv:2604.26752}}
\begin{itemize}
  \item GLM-V 团队多模态智能体基座模型技术报告，参与 GUI Agent 相关后训练工作。
\end{itemize}

%==================== 竞赛 & 荣誉 ====================
\section{竞赛与荣誉}
\entry{程序设计竞赛}{2022 -- 2024}
\begin{itemize}
  \item CCPC 全国邀请赛（广东赛区）\textbf{金牌}；ICPC 南京赛区 \textbf{银牌}、香港赛区 \textbf{银牌}；CCCC 天梯赛个人 \textbf{国家二等奖}（均为 2024）。
\end{itemize}
\vspace{1pt}
\entry{荣誉奖项}{2021 -- 2025}
\begin{itemize}
  \item 上海市优秀毕业生（2025）；国家奖学金（2022）；同济大学电子与信息工程学院启迪奖学金（2023）。
\end{itemize}

\end{document}
